\chapter{Discussion, Limitations and Engineering Significance}

\section{Engineering Significance}
The main engineering significance of Phase 2 is the connection between
machine learning and physical-layer receiver operation. In Phase 1, the
classifier produces a prediction that ends at the classification metric. In
Phase 2, the prediction influences the selected demodulator and therefore
affects the recovered information. This creates a measurable chain from
received waveform to communication reliability.

\section{Why Adaptation is Required}
A fixed receiver is designed for one expected modulation. Such a receiver may
be simple, but it fails when the transmitter changes its format. Adaptive
modulation recognition enables the receiver to support mixed traffic and
potentially select more efficient modulation modes when the channel permits.
The packet-success and throughput results show that high-order modulation is
useful only when the receiver can identify and recover it reliably.

\section{Failure Modes}
The main failure mode is modulation ambiguity at low SNR. Noise causes
different constellations to overlap after normalization, so the minimum
distance decision can select the wrong format. A second failure mode is
residual synchronization error. Even when the modulation is correct, phase,
frequency, or timing error can move symbols across decision boundaries. A
third failure mode occurs when a single bit error causes a packet to be
declared unsuccessful, making packet success stricter than BER.

\section{Limitations}
The present system is a controlled software simulation. The receiver uses a
pilot-assisted channel and carrier estimate, whereas a practical SDR must
estimate these quantities from noisy observations. The current Phase 2
modulation estimator is constellation-based and does not yet invoke the
trained CNN online. The simulation also omits pulse shaping, matched
filtering, timing recovery, forward-error correction, interleaving, and
automatic repeat request.

\section{Reproducibility and Data Integrity}
Every reported result is generated from code in the repository. The random
seed, SNR sweep, number of frames, modulation set, channel model, and output
paths are explicit. The JSON result file is retained alongside the figures,
which provides a machine-readable record for panel verification and future
experiments.

\section{Proposed Phase 3 Extensions}
The next stage should replace the simulated correction with a complete SDR
front-end. A confidence-aware receiver could select a robust fallback
modulation or request retransmission when the classification margin is small.
Forward-error correction could be used to compare uncoded and coded BER.
Further studies should include frequency-selective multipath, Doppler,
oscillator phase noise, nonlinear amplification, and over-the-air recordings.

\section{Summary of New Work}
Compared with Phase 1, the project now includes a transmitter, a configurable
channel, a receiver front-end, adaptive modulation selection, demodulation,
bit recovery, and communication-level evaluation. The new contribution is
therefore a complete adaptive communication experiment rather than only a
change in neural-network architecture.
